\documentclass[letter, 10pt]{article} 
\usepackage[utf8]{inputenc} 
\usepackage[spanish]{babel} 
\usepackage{amsfonts} 
\usepackage{amsmath} 
\usepackage{graphicx} 
\usepackage{url} 
\usepackage{booktabs} 
\usepackage{float} 
\usepackage{pgfplots} 
\pgfplotsset{compat=1.18} \usepackage[top=3cm,bottom=3cm,left=3.5cm,right=3.5cm,footskip=1.5cm,headheight=1.5cm,headsep=.5cm]{geometry}



\begin{document}
\title{Inteligencia Artificial \\ \begin{Large}Estado del arte del problema de localización de instalaciones capacitadas con incompatibilidades entre clientes
\end{Large}}
\author{Nelson Antonio Sepúlveda Valenzuela}
\date{\today}
\maketitle


%--------------------No borrar esta secci\'on--------------------------------%
\section*{Evaluaci\'on}

\begin{tabular}{ll}
Resumen (5\%): & \underline{\hspace{2cm}} \\
Introducci\'on (5\%):  & \underline{\hspace{2cm}} \\
Definici\'on del Problema (10\%):  & \underline{\hspace{2cm}} \\
Estado del Arte (40\%):  & \underline{\hspace{2cm}} \\
Modelo Matem\'atico (20\%): &  \underline{\hspace{2cm}}\\
Conclusiones (15\%): &  \underline{\hspace{2cm}}\\
Bibliograf\'ia (5\%): & \underline{\hspace{2cm}}\\
 &  \\
\textbf{Nota Final (100\%)}:   & \underline{\hspace{2cm}}
\end{tabular}
%---------------------------------------------------------------------------%
\vspace{2cm}

% Resumen %
\begin{abstract}
\noindent Este informe estudia el Multi-Source Capacitated Facility Location Problem with Customer Incompatibilities (MS-CFLP-CI), variante del problema de localización de instalaciones capacitadas. El problema busca decidir qué bodegas abrir y cómo distribuir la demanda de los clientes, considerando costos, capacidades e incompatibilidades. Además, se presenta una representación computacional de las soluciones, una búsqueda local basada en Hill Climbing + Restart y casos de prueba para evaluar su comportamiento. La variante multi-source permite atender clientes desde más de una bodega, aumentando la flexibilidad y complejidad. Finalmente, se exponen conclusiones del estudio bibliográfico y experimental.
\end{abstract}

% Introducción %
\section{Introducción}

\noindent En este informe se estudia el Multi-Source Capacitated Facility Location Problem with Customer Incompatibilities (MS-CFLP-CI), una variante del problema de localización de instalaciones capacitadas aplicada a redes logísticas. El problema consiste en decidir qué bodegas abrir y cómo distribuir la demanda de los clientes entre ellas, considerando costos de apertura, costos de operación, capacidades limitadas e incompatibilidades entre ciertos pares de clientes \cite{klose2005facility,maia2023metaheuristic}.\\

\noindent La motivación de este estudio se relaciona con escenarios logísticos donde no basta con minimizar costos, sino que también se deben cumplir restricciones operacionales. En particular, ciertos clientes o productos pueden no ser compatibles dentro de una misma bodega debido a condiciones de almacenamiento, seguridad, regulación o manejo operacional. Por ello, una solución factible debe satisfacer completamente la demanda de cada cliente, respetar la capacidad disponible en cada bodega y evitar que clientes incompatibles sean atendidos por una misma instalación.\\

\noindent En la variante multi-source, un cliente puede recibir productos desde más de una bodega, lo que entrega mayor flexibilidad para utilizar las capacidades disponibles. Sin embargo, esta característica también aumenta la complejidad del problema, ya que las incompatibilidades deben verificarse en cada bodega donde exista una asignación positiva de demanda \cite{mansini2022kernel,ceschia2024multineighborhood}.\\

\noindent Para abordar el problema, se presenta una representación computacional de las soluciones, junto con una estrategia heurística basada en Hill Climbing + Restart. Este acercamiento permite explorar soluciones vecinas, mejorar progresivamente el costo total y reducir la probabilidad de quedar atrapado en óptimos locales mediante reinicios. Además, se consideran casos de prueba de distinto tamaño para analizar el comportamiento del algoritmo en términos de factibilidad, calidad de solución y desempeño computacional.\\

\noindent El propósito del documento es entregar una base conceptual, formal y experimental para el estudio del MS-CFLP-CI. Para ello, se define el problema y sus componentes principales, se revisa el estado del arte asociado a enfoques exactos, heurísticos, metaheurísticos y matheurísticos, se presenta un modelo matemático, se describe la representación utilizada y se expone el procedimiento de búsqueda local implementado. Finalmente, se analizan los resultados obtenidos en los casos de prueba y se presentan las conclusiones del trabajo.

% Definición del Problema %
\section{Definición del Problema}

\noindent El problema estudiado corresponde al Multi-Source Capacitated Facility Location Problem with Customer Incompatibilities (MS-CFLP-CI), una variante del problema de localización de instalaciones capacitadas con incompatibilidades entre clientes. En términos generales, el problema consiste en decidir qué bodegas, instalaciones o centros de distribución deben abrirse y cómo debe distribuirse la demanda de los clientes entre ellas. Cada bodega posee una capacidad máxima y un costo fijo de apertura, mientras que atender a un cliente desde una bodega genera un costo operacional asociado. El objetivo es encontrar una solución de costo mínimo, considerando tanto los costos de apertura como los costos de operación necesarios para satisfacer completamente la demanda de todos los clientes \cite{klose2005facility,maia2023metaheuristic,ceschia2024multineighborhood}.\\

\noindent La característica principal del MS-CFLP-CI es que un cliente puede recibir su demanda desde más de una bodega. Es decir, la demanda de un cliente puede dividirse entre distintas instalaciones, siempre que la suma total recibida satisfaga completamente su requerimiento. Esta condición diferencia a la variante multi-source de la variante single-source, donde cada cliente debe ser atendido por una única instalación. La posibilidad de dividir la demanda entrega mayor flexibilidad al modelo, ya que permite aprovechar mejor las capacidades disponibles y combinar alternativas de menor costo \cite{mansini2022kernel,ceschia2024multineighborhood}.\\

\noindent Además de las restricciones clásicas de capacidad y demanda, el MS-CFLP-CI incorpora incompatibilidades entre clientes. Estas incompatibilidades se representan mediante pares de clientes que no pueden ser atendidos por una misma bodega. En la práctica, esta situación puede aparecer cuando ciertos productos, servicios o requerimientos no deben compartir una misma instalación por razones sanitarias, técnicas, operativas o de seguridad. Por lo tanto, una solución factible debe asegurar que cada cliente reciba toda su demanda, que ninguna bodega exceda su capacidad y que ningún par de clientes incompatibles sea servido desde la misma instalación \cite{maia2023metaheuristic,mansini2022kernel,ceschia2024multineighborhood}.\\

\noindent Las decisiones principales del problema se pueden describir en tres grupos. Primero, se debe determinar qué bodegas serán abiertas, lo que implica asumir sus respectivos costos fijos. Segundo, se debe decidir cuánta demanda de cada cliente será enviada desde cada bodega abierta. Tercero, se debe verificar que las asignaciones realizadas no generen conflictos de incompatibilidad. En la variante MS-CFLP-CI, esta última verificación es especialmente importante, ya que un cliente queda asociado a cada bodega desde la cual recibe una cantidad positiva de productos \cite{mansini2022kernel,ceschia2024multineighborhood}.\\

\noindent Desde el punto de vista computacional, el MS-CFLP-CI es un problema de alta dificultad, ya que hereda la complejidad del CFLP clásico, reconocido como un problema NP-Hard, e incorpora restricciones adicionales de incompatibilidad entre clientes. Además, el tamaño del espacio de búsqueda crece rápidamente: para $m$ bodegas existen $2^m$ posibles combinaciones de apertura, y para $n$ clientes la asignación multi-source puede considerar múltiples subconjuntos de bodegas para cada cliente. Esto genera un espacio combinatorio muy grande, especialmente en instancias con muchas instalaciones, clientes e incompatibilidades \cite{klose2005facility,maia2023metaheuristic,ceschia2024multineighborhood}.\\

\noindent El MS-CFLP-CI también puede relacionarse con otras variantes de problemas de localización, como el Uncapacitated Facility Location Problem (UFLP), donde no se consideran límites de capacidad; el $p$-Median Problem, que busca ubicar un número fijo de instalaciones minimizando distancias o costos de asignación; y el Robust Facility Location Problem, que incorpora incertidumbre en parámetros como demanda, costos o disponibilidad. A diferencia de estas variantes, el MS-CFLP-CI combina capacidades limitadas, asignación divisible de demanda e incompatibilidades entre clientes, por lo que requiere enfoques de solución más especializados, incluyendo modelos matemáticos, heurísticas, metaheurísticas y matheurísticas \cite{fischetti2016benders,maia2023metaheuristic,mansini2022kernel,ceschia2024multineighborhood}.

% Estado del Arte %
\section{Estado del Arte}

\noindent El problema de localización de instalaciones es una línea clásica de estudio dentro de la investigación de operaciones, la optimización combinatoria y el diseño de redes logísticas. Su propósito general es determinar en qué ubicaciones abrir instalaciones y cómo asignar clientes a ellas, de modo que la demanda sea satisfecha al menor costo posible. Dentro de esta familia de problemas, el Capacitated Facility Location Problem (CFLP) incorpora capacidades máximas en las instalaciones, lo que permite representar sistemas donde las bodegas, plantas o centros de distribución no pueden atender una cantidad ilimitada de demanda \cite{klose2005facility}.\\

\noindent Históricamente, el CFLP ha sido abordado mediante modelos de programación matemática, especialmente formulaciones de programación lineal entera mixta. Estas representaciones consideran decisiones de apertura de instalaciones, asignación de clientes y uso de capacidad, permitiendo obtener soluciones óptimas en instancias pequeñas o medianas. Entre los enfoques exactos más relevantes se encuentran métodos de ramificación, planos de corte y descomposición de Benders. Sin embargo, la dificultad computacional aumenta significativamente cuando se incorporan restricciones adicionales, como capacidades ajustadas, asignación única o incompatibilidades entre clientes \cite{fischetti2016benders,klose2005facility}.\\

\noindent El Capacitated Facility Location Problem with Customer Incompatibilities (CFLP-CI) surge como una extensión del CFLP clásico al incorporar restricciones de incompatibilidad entre clientes. En esta variante, ciertos pares de clientes no pueden ser atendidos por una misma instalación, lo que permite modelar situaciones donde algunos productos, servicios o requerimientos no deben coexistir en una misma bodega por razones sanitarias, técnicas, operacionales o de seguridad. Esta característica modifica la estructura del problema, ya que la factibilidad de una solución no depende únicamente de satisfacer demanda y respetar capacidades, sino también de evitar conflictos entre clientes asignados a una misma instalación \cite{maia2023metaheuristic,ceschia2024multineighborhood}.\\

\noindent Desde el punto de vista del modelamiento, el CFLP-CI suele representarse mediante formulaciones de programación entera mixta que combinan variables de apertura de instalaciones, variables de flujo o asignación de demanda y variables binarias auxiliares para indicar si un cliente es atendido por una determinada bodega. Además, las incompatibilidades pueden representarse como un grafo de conflictos, donde los nodos corresponden a clientes y las aristas indican pares incompatibles. Esta representación es efectiva porque permite traducir las incompatibilidades en restricciones que impiden que dos clientes conectados por una arista sean servidos por la misma instalación \cite{maia2023metaheuristic,mansini2022kernel,ceschia2024multineighborhood}.\\

\noindent Una de las primeras aproximaciones específicas al CFLP-CI fue desarrollada por Maia et al., quienes propusieron y compararon distintas técnicas metaheurísticas para resolver instancias del problema. Entre los enfoques considerados se encuentran métodos evolutivos, búsqueda local, procedimientos greedy y estrategias híbridas basadas en minería de datos. Sus resultados muestran que las metaheurísticas son una alternativa adecuada para abordar instancias de mayor tamaño, donde los métodos exactos pueden requerir tiempos computacionales elevados o no alcanzar soluciones óptimas dentro de límites razonables \cite{maia2023metaheuristic}.\\

\noindent Posteriormente, Mansini y Zanotti abordaron problemas de localización con incompatibilidades desde una perspectiva matheurística, considerando variantes single-source y multi-source. Su propuesta se basa en un enfoque de Two-Phase Kernel Search, que combina programación matemática con una estrategia de reducción del espacio de búsqueda. Este tipo de método resulta relevante porque conserva parte de la estructura formal del modelo, pero evita explorar todas las variables posibles, concentrándose en subconjuntos prometedores para obtener soluciones de buena calidad en menor tiempo \cite{mansini2022kernel}.\\

\noindent La variante de mayor interés para este informe es el Multi-Source Capacitated Facility Location Problem with Customer Incompatibilities (MS-CFLP-CI), donde un cliente puede recibir demanda desde más de una bodega. Esta flexibilidad permite aprovechar mejor la capacidad disponible y reducir costos de operación, pero también aumenta la complejidad, porque las incompatibilidades deben verificarse en cada bodega donde un cliente recibe una cantidad positiva de productos. En esta línea, el trabajo de Ceschia y Schaerf representa uno de los avances más relevantes, al proponer un algoritmo de Simulated Annealing multi-vecindario diseñado específicamente para esta variante. En síntesis, la tendencia actual apunta al uso de metaheurísticas especializadas y matheurísticas, ya que los métodos exactos siguen siendo útiles para formular el problema, obtener cotas y resolver instancias pequeñas, mientras que los enfoques híbridos permiten abordar instancias grandes y altamente restringidas con mayor eficiencia computacional \cite{ceschia2024multineighborhood,mansini2022kernel,fischetti2016benders}.

% Modelo Matemático %
\section{Modelo matemático}

\noindent En esta sección se presenta un modelo matemático para el Multi-Source Capacitated Facility Location Problem with Customer Incompatibilities (MS-CFLP-CI). Este modelo representa las decisiones de apertura de bodegas, asignación de demanda de clientes y control de incompatibilidades. La formulación corresponde a un modelo de programación lineal entera mixta, ya que combina variables binarias con variables de cantidad o flujo \cite{maia2023metaheuristic,mansini2022kernel,ceschia2024multineighborhood}.

\subsection{Conjuntos}

\noindent Sea $I$ el conjunto de bodegas candidatas y sea $J$ el conjunto de clientes. Además, se define $P$ como el conjunto de pares de clientes incompatibles. Si $(a,b) \in P$, entonces los clientes $a$ y $b$ no pueden ser atendidos por una misma bodega.

\subsection{Parámetros}

\noindent Los parámetros considerados son los siguientes:

\begin{itemize}
    \item $C_i$: capacidad máxima de la bodega $i$.
    \item $FC_i$: costo fijo de apertura de la bodega $i$.
    \item $G_j$: demanda total del cliente $j$.
    \item $SC_{ji}$: costo unitario de operación por satisfacer demanda del cliente $j$ desde la bodega $i$.
\end{itemize}

\subsection{Variables de decisión}

\noindent El modelo utiliza las siguientes variables:

\begin{itemize}
    \item $y_i \in \{0,1\}$: toma valor $1$ si la bodega $i$ es abierta, y $0$ en caso contrario.
    \item $x_{ji} \in \mathbb{Z}_{\geq 0}$: cantidad de demanda del cliente $j$ satisfecha desde la bodega $i$.
    \item $z_{ji} \in \{0,1\}$: toma valor $1$ si el cliente $j$ recibe una cantidad positiva de productos desde la bodega $i$, y $0$ en caso contrario.
\end{itemize}

\noindent La variable $x_{ji}$ permite representar la característica principal del MS-CFLP-CI, ya que un cliente puede recibir demanda desde más de una bodega. Por otra parte, la variable $z_{ji}$ permite identificar si existe una asignación efectiva entre un cliente y una bodega, lo que resulta necesario para controlar las incompatibilidades entre clientes.

\subsection{Función objetivo}

\noindent El objetivo es minimizar el costo total de la solución, compuesto por los costos fijos de apertura de bodegas y los costos operacionales asociados al envío de demanda.

\begin{equation}
\min
\left\{
\sum_{i \in I} FC_i y_i
+
\sum_{j \in J} \sum_{i \in I} SC_{ji} x_{ji}
\right\}
\end{equation}

\noindent El primer término representa el costo de abrir las bodegas utilizadas, mientras que el segundo término representa el costo de satisfacer la demanda de los clientes desde dichas bodegas.

\subsection{Restricciones}

\noindent Cada cliente debe tener su demanda completamente satisfecha. La demanda puede ser cubierta por una o varias bodegas.

\begin{equation}
\sum_{i \in I} x_{ji} = G_j \quad \forall j \in J
\end{equation}

\noindent Cada bodega solo puede operar hasta su capacidad máxima. Además, si una bodega no está abierta, no puede enviar productos a ningún cliente.

\begin{equation}
\sum_{j \in J} x_{ji} \leq C_i y_i \quad \forall i \in I
\end{equation}

\newpage

\noindent Las siguientes restricciones relacionan la cantidad enviada con la variable de asignación. Si el cliente $j$ no está asignado a la bodega $i$, entonces no puede recibir demanda desde ella. A su vez, si $z_{ji}$ toma valor $1$, se exige que exista una cantidad positiva enviada desde la bodega $i$ hacia el cliente $j$.

\begin{equation}
x_{ji} \leq G_j z_{ji} \quad \forall j \in J, \forall i \in I
\end{equation}

\begin{equation}
x_{ji} \geq z_{ji} \quad \forall j \in J, \forall i \in I
\end{equation}

\noindent También se exige que un cliente solo pueda ser atendido por una bodega que esté abierta.

\begin{equation}
z_{ji} \leq y_i \quad \forall j \in J, \forall i \in I
\end{equation}

\noindent Finalmente, se incorporan las incompatibilidades entre clientes. Si dos clientes son incompatibles, no pueden ser atendidos por una misma bodega.

\begin{equation}
z_{ai} + z_{bi} \leq 1 \quad \forall (a,b) \in P, \forall i \in I
\end{equation}

\noindent El dominio de las variables queda definido como:

\begin{equation}
y_i \in \{0,1\} \quad \forall i \in I
\end{equation}

\begin{equation}
z_{ji} \in \{0,1\} \quad \forall j \in J, \forall i \in I
\end{equation}

\begin{equation}
x_{ji} \in \mathbb{Z}_{\geq 0} \quad \forall j \in J, \forall i \in I
\end{equation}


% Representación %
\section{Representación}

\noindent Para representar una solución del MS-CFLP-CI se utiliza una matriz de envíos, cuyas filas corresponden a los clientes y cuyas columnas corresponden a las bodegas candidatas. Cada elemento \(x_{ji}\) indica la cantidad de demanda del cliente \(j\) que es satisfecha desde la bodega \(i\).\\

\noindent Esta representación es apropiada para la variante multi-source, ya que permite que un mismo cliente sea atendido por más de una bodega. En consecuencia, una fila puede contener varios valores positivos, indicando que la demanda del cliente se divide entre distintas instalaciones. Esta característica permite representar soluciones que no serían posibles en una variante single-source, donde cada cliente solo podría ser asignado a una bodega.\\

\noindent Por ejemplo, una solución con tres clientes y dos bodegas puede representarse como:

\[
X =
\begin{pmatrix}
5 & 0 \\
2 & 4 \\
0 & 6
\end{pmatrix}
\]

\noindent En esta matriz, el primer cliente es atendido completamente por la primera bodega, el segundo cliente recibe productos desde ambas bodegas y el tercer cliente es atendido solo por la segunda bodega. Por lo tanto, la representación permite expresar directamente el carácter multi-source del problema.\\

\noindent A partir de la matriz se puede evaluar la factibilidad de una solución. La demanda de cada cliente se obtiene sumando los valores de su fila, mientras que la capacidad utilizada por cada bodega se calcula sumando los valores de su columna. 

\noindent Además, la apertura de una bodega se determina a partir de la misma matriz. Si una columna posee al menos un valor positivo, entonces la bodega correspondiente se considera abierta y se incorpora su costo fijo en la función objetivo. Si todos los valores de la columna son cero, la bodega no participa en la solución.\\

\noindent Las incompatibilidades entre clientes también se verifican mediante la matriz de envíos. Para cada par de clientes incompatibles \((a,b)\), no debe existir una bodega \(i\) desde la cual ambos reciban una cantidad positiva. 

\noindent La representación mantiene una relación directa con las soluciones del problema, ya que cada matriz factible define una asignación válida de demanda y cada asignación válida puede expresarse mediante una matriz de este tipo. Además, no excluye soluciones óptimas, pues permite representar cualquier distribución de demanda entre bodegas que respete las restricciones del MS-CFLP-CI.\\

\noindent Esta estructura también se complementa adecuadamente con el método Hill Climbing + Restart, ya que los movimientos pueden definirse como modificaciones sobre los valores de la matriz, tales como reasignar demanda entre bodegas, abrir o cerrar instalaciones utilizadas, o redistribuir parcialmente la demanda de un cliente. De esta manera, la representación permite explorar soluciones factibles e infactibles durante la búsqueda, manteniendo una forma simple e informativa para evaluar costos, capacidades e incompatibilidades.

% Descripción del algoritmo %
\section{Descripción del Algoritmo}

\noindent Para resolver el MS-CFLP-CI se propone un algoritmo basado en Hill Climbing con estrategia de Best Improvement y múltiples reinicios. Esta elección se justifica porque el problema posee un espacio de búsqueda amplio y restringido por condiciones de capacidad, demanda e incompatibilidad entre clientes. En este contexto, la búsqueda local permite mejorar progresivamente una solución factible, mientras que los reinicios incorporan diversificación y reducen la dependencia respecto de una única solución inicial.\\

\noindent El algoritmo comienza con la generación de una solución inicial factible mediante una construcción greedy aleatorizada. En este proceso, los clientes se asignan a bodegas que respeten la capacidad disponible y las restricciones de incompatibilidad. Normalmente se priorizan asignaciones de menor costo incremental, pero se incorpora un componente aleatorio que permite seleccionar alternativas distintas dentro de un conjunto de candidatos factibles. Este parámetro de aleatoriedad favorece que cada reinicio explore una región diferente del espacio de soluciones.\\

\noindent Una vez construida la solución inicial, se aplica una búsqueda local mediante Hill Climbing Best Improvement. En cada iteración se genera un vecindario a partir de la solución actual. Los movimientos utilizados modifican la matriz de envíos, por ejemplo, reasignando la demanda de un cliente hacia otra bodega factible, redistribuyendo parcialmente su demanda entre distintas bodegas o eliminando envíos que permitan cerrar una instalación cuando esta deja de ser utilizada. Estos operadores son coherentes con la representación elegida, ya que actúan directamente sobre las cantidades enviadas y permiten modificar las decisiones de apertura y asignación.\\

\noindent La estrategia Best Improvement evalúa todos los vecinos factibles generados en una iteración y selecciona aquel que produzca la mayor reducción en la función objetivo. Sea \(S\) la solución actual y \(N(S)\) su vecindario factible, se selecciona una solución \(S'\) tal que:

\[
S' = \arg\min_{S_k \in N(S)} FO(S_k),
\]

\noindent siempre que se cumpla:

\[
FO(S') < FO(S).
\]

\noindent Si ningún vecino factible mejora el valor de la función objetivo, la solución actual se considera un óptimo local y finaliza la búsqueda local de ese reinicio. Aunque esta estrategia puede requerir más tiempo por iteración que First Improvement, permite elegir el mejor movimiento disponible en cada paso, favoreciendo una intensificación más ordenada de la búsqueda.\\

\noindent Para reducir el riesgo de quedar limitado a un óptimo local, se incorpora una estrategia de restart. En cada reinicio se genera una nueva solución inicial factible utilizando una semilla distinta y luego se ejecuta nuevamente la búsqueda local. Al finalizar cada reinicio, el óptimo local obtenido se compara con la mejor solución global encontrada hasta el momento. Si su costo es menor, esta solución pasa a ser la mejor solución conocida.\\

\noindent El algoritmo finaliza cuando se alcanza el número máximo de reinicios, el límite de iteraciones o el tiempo máximo definido. Estos parámetros controlan el esfuerzo computacional del método: un mayor número de reinicios aumenta la diversificación, mientras que un mayor número de iteraciones permite intensificar la búsqueda en cada región explorada. En conjunto, el algoritmo propuesto mantiene un equilibrio entre intensificación y diversificación, lo que resulta apropiado para el MS-CFLP-CI debido a la complejidad generada por las capacidades limitadas, los costos de apertura y las incompatibilidades entre clientes.

% Experimentos %
\section{Experimentos}

\noindent La experimentación tuvo como objetivo evaluar el comportamiento del algoritmo propuesto para el MS-CFLP-CI en términos de calidad de solución, factibilidad y tiempo de ejecución. Para ello, se utilizó una implementación en C++ basada en Hill Climbing con estrategia Best Improvement y reinicios. El algoritmo trabaja sobre una representación mediante matriz de envíos \(x_{ji}\), lo que permite modificar las asignaciones entre clientes y bodegas manteniendo consistencia con la variante multi-source del problema.\\

\noindent Las ejecuciones se realizaron en un computador Acer Predator Triton 300 SE, modelo Predator PT314-51s, con procesador 11th Gen Intel(R) Core(TM) i7-11375H @ 3.30GHz, \(16\) GB de memoria RAM y sistema operativo de \(64\) bits. Para mantener una comparación consistente, todas las instancias fueron ejecutadas utilizando el mismo código fuente, compilador y configuración experimental. Como indicadores de evaluación se consideraron el valor final de la función objetivo, la factibilidad de la solución obtenida, el tiempo de ejecución y la mejora alcanzada respecto de una solución inicial greedy.\\

\noindent El procedimiento experimental consideró tres configuraciones principales. Primero, se evaluó una construcción greedy determinística, equivalente a ejecutar el algoritmo con \(\epsilon = 0\), un solo reinicio y sin iteraciones de búsqueda local. Esta configuración se utilizó como línea base. Luego, se ejecutó Hill Climbing Best Improvement desde una solución inicial determinística, con \(\epsilon = 0\), un reinicio y un número fijo de iteraciones. Finalmente, se evaluó la versión estocástica con reinicios, donde cada reinicio genera una solución inicial mediante una regla \(\epsilon\)-greedy.\\

\noindent Para todas las configuraciones principales se utilizaron cinco semillas: \(101\), \(202\), \(303\), \(404\) y \(505\). La versión con \(\epsilon = 0\) es determinística, ya que bajo las mismas condiciones construye la misma solución inicial y aplica el mismo criterio de mejora. En cambio, la versión con \(\epsilon > 0\) y reinicios es estocástica, debido a que distintas semillas pueden generar soluciones iniciales diferentes. Esta diferencia permite analizar el efecto de la diversificación en la calidad de las soluciones obtenidas.\\

\noindent Para estudiar el efecto del parámetro \(\epsilon\), se utilizó la instancia wlp01.in, compuesta por \(50\) bodegas y \(115\) clientes. Esta instancia fue seleccionada por tener un tamaño intermedio pequeño, lo que permite ejecutar varias semillas en tiempos razonables. Se probaron los valores \(\epsilon \in \{0, 0.10, 0.20, 0.40\}\), con \(5\) reinicios, \(30\) iteraciones máximas por reinicio y las cinco semillas definidas. A partir de estos resultados se seleccionó \(\epsilon = 0.20\), ya que entregó el menor costo promedio.\\

\noindent Posteriormente, se estudió el desempeño del algoritmo en instancias de distinto tamaño. Se utilizaron las instancias toy.in, wlp01.in, wlp02.in, wlp03.in, wlp04.in, wlp05.in, wlp06.in y wlp07.in, cubriendo más de cinco casos de prueba. Para toy.in, wlp01.in, wlp02.in y wlp03.in se utilizó la configuración \(\epsilon = 0.20\), \(5\) reinicios y \(30\) iteraciones por reinicio. Para las instancias más grandes, desde wlp04.in hasta wlp07.in, se redujo la configuración a \(3\) reinicios y \(20\) iteraciones por reinicio, con el objetivo de mantener tiempos de ejecución razonables debido al crecimiento del vecindario.\\

\noindent Además, se realizó un análisis de convergencia sobre la instancia wlp03.in variando el número máximo de iteraciones por reinicio en \(\{0, 5, 10, 20, 30\}\). Este experimento permitió observar cómo cambia la calidad de la solución respecto del avance de la búsqueda y del tiempo de ejecución.\\

\noindent El criterio de término del algoritmo estuvo dado por el número máximo de reinicios y el número máximo de iteraciones por reinicio. Además, se consideró un límite máximo de tiempo por ejecución para evitar corridas excesivamente largas. En cada ejecución se registró la mejor solución encontrada, su costo total, su factibilidad y el tiempo requerido. De esta forma, el diseño experimental permite comparar la línea base greedy con la búsqueda local determinística, la versión estocástica con reinicios y el comportamiento de convergencia del método.

% Resultados %
\section{Resultados}

\noindent Los resultados fueron obtenidos a partir de cinco semillas distintas: \(101\), \(202\), \(303\), \(404\) y \(505\). En todos los casos, las soluciones reportadas fueron factibles, ya que el programa verifica el cumplimiento de demanda, capacidad e incompatibilidades antes de entregar el resultado final.\\

\noindent La Tabla~\ref{tab:epsilon} muestra el efecto del parámetro \(\epsilon\) sobre la instancia wlp01.in. Se observa que la configuración determinística, con \(\epsilon = 0\), obtiene siempre el mismo costo, lo que confirma la ausencia de diversificación. Al incorporar aleatoriedad, el costo promedio disminuye considerablemente. El mejor promedio se obtiene con \(\epsilon = 0.20\), aunque \(\epsilon = 0.10\) alcanza la mejor corrida individual. Por esta razón, se selecciona \(\epsilon = 0.20\) para los experimentos restantes, al presentar el mejor comportamiento promedio.

\begin{table}[H]
\centering
\caption{Ajuste de \(\epsilon\) en wlp01.in con 5 reinicios, 30 iteraciones y 5 semillas.}
\label{tab:epsilon}
\begin{tabular}{rrrrr}
\toprule
\(\epsilon\) & Costo promedio & Mejor costo & Peor costo & Tiempo promedio (s) \\
\midrule
0.00 & 47249.0 & 47249 & 47249 & 0.035 \\
0.10 & 41704.8 & 39285 & 42867 & 0.038 \\
0.20 & 40666.0 & 39540 & 41413 & 0.040 \\
0.40 & 41655.4 & 41172 & 42273 & 0.043 \\
\bottomrule
\end{tabular}
\end{table}

\noindent La Tabla~\ref{tab:resultados-generales} compara la construcción greedy, el Hill Climbing determinístico y la versión con reinicios. Los valores corresponden al promedio de las cinco semillas. En todas las instancias, la búsqueda local mejora la solución greedy inicial, y la versión con reinicios obtiene los menores costos. Esto evidencia que la combinación entre una construcción \(\epsilon\)-greedy y reinicios permite explorar regiones distintas del espacio de búsqueda.

\begin{table}[H]
\centering
\caption{Resultados promedio sobre 5 semillas.}
\label{tab:resultados-generales}
\begin{tabular}{lrrrrrrr}
\toprule
Instancia & Bod. & Cli. & Greedy & HC det. & HC restart & Mejor & Mejora \% \\
\midrule
\texttt{toy.in}   & 4   & 10   & 8586.0   & 7989.0   & 6801.0   & 6757   & 20.79 \\
\texttt{wlp01.in} & 50  & 115  & 53661.0  & 47249.0  & 40666.0  & 39540  & 24.22 \\
\texttt{wlp02.in} & 100 & 253  & 117872.0 & 103323.0 & 90242.2  & 85181  & 23.44 \\
\texttt{wlp03.in} & 150 & 345  & 150386.0 & 135650.0 & 116414.4 & 115082 & 22.59 \\
\texttt{wlp04.in} & 200 & 479  & 192426.0 & 183794.0 & 166165.4 & 160350 & 13.65 \\
\texttt{wlp05.in} & 250 & 601  & 238818.0 & 229106.0 & 208611.0 & 204281 & 12.65 \\
\texttt{wlp06.in} & 300 & 705  & 273461.0 & 263636.0 & 233144.4 & 231395 & 14.74 \\
\texttt{wlp07.in} & 400 & 1012 & 382809.0 & 371740.0 & 340005.2 & 335820 & 11.18 \\
\bottomrule
\end{tabular}
\end{table}

\noindent En las instancias pequeñas y medianas, la mejora respecto de la solución greedy se mantiene cercana o superior al \(20\%\). En las instancias grandes, la mejora porcentual disminuye, pero sigue siendo consistente. Esto se explica porque, al crecer el número de clientes y bodegas, también aumenta el tamaño del vecindario y se vuelve más difícil explorar una proporción significativa del espacio de soluciones con la misma cantidad de iteraciones.

\noindent La Tabla~\ref{tab:tiempos} presenta el tiempo promedio de la versión con reinicios. Como era esperable, el costo computacional aumenta con el tamaño de la instancia, debido a que en cada iteración se evalúan más movimientos posibles entre clientes y bodegas. Aun así, los tiempos se mantienen razonables para una primera implementación heurística.

\begin{table}[H]
\centering
\caption{Tiempo promedio de HC con reinicios sobre 5 semillas.}
\label{tab:tiempos}
\begin{tabular}{lrr}
\toprule
Instancia & Tiempo promedio (s) & Mejor costo obtenido \\
\midrule
\texttt{toy.in}   & 0.003  & 6757 \\
\texttt{wlp01.in} & 0.040  & 39540 \\
\texttt{wlp02.in} & 0.250  & 85181 \\
\texttt{wlp03.in} & 0.916  & 115082 \\
\texttt{wlp04.in} & 1.515  & 160350 \\
\texttt{wlp05.in} & 3.373  & 204281 \\
\texttt{wlp06.in} & 5.682  & 231395 \\
\texttt{wlp07.in} & 17.358 & 335820 \\
\bottomrule
\end{tabular}
\end{table}

\noindent Para analizar la convergencia, se utilizó la instancia wlp03.in variando el número máximo de iteraciones por reinicio. La Tabla~\ref{tab:convergencia} muestra que, a medida que aumenta el número de iteraciones, el costo promedio disminuye. Las mejoras son progresivas, aunque cada aumento de iteraciones implica un mayor esfuerzo computacional. Este comportamiento es coherente con una búsqueda local, ya que las primeras iteraciones suelen encontrar mejoras rápidas y luego el algoritmo se aproxima gradualmente a óptimos locales.

\begin{table}[H]
\centering
\caption{Análisis de convergencia en wlp03.in con \(\epsilon = 0.20\), 5 reinicios y 5 semillas.}
\label{tab:convergencia}
\begin{tabular}{rrrrr}
\toprule
Iteraciones por reinicio & Costo promedio & Mejor costo & Tiempo promedio (s) & Mejora \% \\
\midrule
0  & 130996.6 & 129455 & 0.673 & 0.00 \\
5  & 127242.2 & 125954 & 0.758 & 2.87 \\
10 & 124418.0 & 123166 & 0.683 & 5.02 \\
20 & 119984.2 & 118512 & 0.880 & 8.41 \\
30 & 116414.4 & 115082 & 0.900 & 11.13 \\
\bottomrule
\end{tabular}
\end{table}


% Conclusiones %
\section{Conclusiones}

\noindent Se concluye que el Multi-Source Capacitated Facility Location Problem with Customer Incompatibilities (MS-CFLP-CI) corresponde a una variante relevante y compleja del problema de localización de instalaciones capacitadas. Su principal característica es permitir que la demanda de un cliente sea satisfecha por más de una bodega, incorporando además restricciones de incompatibilidad entre clientes. Esto permite representar escenarios logísticos más realistas, donde no solo se deben minimizar costos de apertura y operación, sino también respetar capacidades e impedir que ciertos pares de clientes sean atendidos por una misma instalación \cite{maia2023metaheuristic,ceschia2024multineighborhood}.\\

\noindent El estudio realizado permite distinguir la importancia de las variantes del problema. El CFLP clásico, el CFLP-CI, la variante single-source y la variante multi-source presentan diferencias relevantes en la forma de representar y resolver las asignaciones. En particular, la variante multi-source entrega mayor flexibilidad, pero también aumenta la complejidad, ya que las incompatibilidades deben verificarse en cada bodega donde exista asignación positiva de demanda.\\

\noindent La propuesta implementada, basada en una representación mediante matriz de envíos y en un algoritmo Hill Climbing con estrategia Best Improvement y reinicios, resultó adecuada para abordar el problema. La matriz de envíos permitió representar soluciones compatibles con la naturaleza multi-source, mientras que los operadores de búsqueda local permitieron modificar asignaciones, redistribuir demanda y mejorar progresivamente el costo total. Además, la construcción \(\epsilon\)-greedy permitió generar soluciones iniciales diversas sin abandonar completamente el criterio de bajo costo.\\

\noindent Los resultados experimentales justifican la relevancia de la propuesta. La solución greedy inicial permitió obtener soluciones factibles de manera rápida, pero su carácter determinístico limitó la exploración del espacio de búsqueda. La aplicación de Hill Climbing determinístico redujo el costo inicial, y la incorporación de reinicios estocásticos permitió obtener soluciones de menor costo en todas las instancias evaluadas. En las instancias pequeñas y medianas, la mejora promedio respecto de la solución greedy estuvo aproximadamente entre \(20.79\%\) y \(24.22\%\). En las instancias grandes, las mejoras fueron menores, entre \(11.18\%\) y \(14.74\%\), pero se mantuvieron consistentes pese al aumento del tamaño del vecindario.\\

\noindent También se concluye que el parámetro \(\epsilon\) influye directamente en la calidad de las soluciones obtenidas. Un valor nulo elimina la diversificación y genera siempre la misma solución inicial, mientras que un valor demasiado alto puede introducir aleatoriedad excesiva y deteriorar la calidad de las soluciones. En los experimentos realizados sobre wlp01.in, el valor \(\epsilon = 0.20\) obtuvo el mejor desempeño promedio, lo que evidencia un equilibrio adecuado entre intensificación greedy y diversificación estocástica.\\

\noindent El análisis de convergencia permitió observar que aumentar el número de iteraciones por reinicio mejora la calidad promedio de las soluciones, aunque con un mayor costo computacional. En la instancia wlp03.in, el costo promedio disminuyó progresivamente al aumentar las iteraciones, lo que confirma que la búsqueda local continúa encontrando mejoras durante el proceso. Sin embargo, el crecimiento del tiempo de ejecución en instancias grandes muestra que la exploración completa del vecindario puede volverse costosa.\\

\noindent Respecto del estado del arte, los métodos exactos, como la programación lineal entera mixta o la descomposición de Benders, son relevantes para formular el problema y obtener soluciones óptimas en instancias pequeñas o medianas, pero pueden presentar dificultades en instancias grandes o con restricciones adicionales \cite{fischetti2016benders}. Por ello, los enfoques heurísticos, metaheurísticos y matheurísticos resultan especialmente importantes para este tipo de problema, ya que permiten obtener soluciones de buena calidad en tiempos computacionales razonables \cite{maia2023metaheuristic,mansini2022kernel}.\\

\noindent Como trabajo futuro, sería relevante incorporar vecindarios más selectivos y estrategias metaheurísticas más robustas, como Simulated Annealing, Tabu Search o enfoques híbridos. También sería útil comparar distintas metaheurísticas bajo las mismas instancias, tiempos de ejecución y criterios de evaluación.

% Uso LLMs %
\section{Uso LLMs}

\noindent Para el desarrollo de este trabajo se utilizó una herramienta de LLM/IA como apoyo en tareas de revisión y mejora del texto.

\begin{enumerate}
    \item \textbf{Herramienta utilizada:} ChatGPT (GPT-5.5 Thinking).

    \item \textbf{Propósito de uso:} La herramienta fue utilizada principalmente para corregir faltas ortográficas, mejorar la redacción, revisar gramática, ordenar ideas y sugerir una formulación más clara de algunos párrafos del informe.

    \item \textbf{Prompts relevantes:} Se realizaron consultas orientadas a mejorar la calidad del texto, tales como: ``corrige la ortografía y redacción de este texto'', ``mejora la claridad de este párrafo'' y ``redacta esta sección en tercera persona y con estilo académico''.

    \item \textbf{Nivel de edición:} Las respuestas generadas fueron revisadas, adaptadas y modificadas antes de ser incorporadas al informe. No se copiaron sin revisión, ya que el contenido fue ajustado según los requerimientos del trabajo y contrastado con las fuentes bibliográficas utilizadas.

    \item \textbf{Validación:} La información técnica fue verificada utilizando papers académicos, el material entregado por el curso y juicio propio. La herramienta se utilizó como apoyo de redacción, no como única fuente de información.
\end{enumerate}

% Referencias %
\bibliographystyle{plain}
\bibliography{Referencias}

\end{document} 
